\documentclass[12pt]{article}
\usepackage[utf8]{inputenc}
\usepackage{amsmath, amssymb}
\usepackage{geometry}
\geometry{margin=1in}

\title{he Specific-Factors Model}
\author{International Economic Relations}
\date{}

\begin{document}
%\maketitle

\section*{Problem (16 points)}

Canada produces two goods: Automobiles ($A$) and Wheat ($W$).  
Labor is mobile between the two sectors, while capital is specific to automobiles and land is specific to wheat.

\[
Q_A = K^{1/2}L_A^{1/2}, 
\qquad 
Q_W = T^{1/2}L_W^{1/2},
\qquad 
L_A + L_W = L.
\]

Endowments:
\[
K = 49, \quad T = 81, \quad L = 120.
\]

Prices: initially $P_A = 4$, $P_W = 2$. Later, $P_A$ rises to $6$ while $P_W$ remains $2$.  

\bigskip

\noindent 1. Derive the marginal products of labor in each sector.  

\bigskip

\noindent 2. Write the wage equalization condition and show the equation that determines the allocation of labor between the two sectors.  

\bigskip

\noindent 3. Solve for the equilibrium allocation of labor ($L_A$, $L_W$) and the common wage when $P_A = 4$.  

\bigskip

\noindent 4. Recalculate the equilibrium labor allocation and wage when $P_A = 6$.  

\bigskip



\noindent 5. Suppose the world price of automobiles is higher than Canada’s initial autarky price.  
Explain, using the logic of the specific-factors model, how opening up to trade would affect:  

\begin{itemize}
    \item (a) Labor allocation between the two sectors,  
    \item (b) The real return to capital, land, and labor, and  
    \item (c) The pattern of trade (which good Canada exports and imports).  
\end{itemize}

\end{document}
